 \paragraph*{04.05.08.05 Ressourcenverbrauch evaluieren und erforderliche Korrekturmaßnahmen ergreifen}
\kcilabel{para:04.05.08.05_Ressourcenverbrauch}{04.05.08.05}
\label{para:04.05.08.05_Ressourcenverbrauch}

% Task
\par Im Projekt \gls{r2h} oblag mir als \gls{pl} die kontinuierliche Überwachung des Ressourcenverbrauchs -- insbesondere hinsichtlich Zeit, Personal und finanzieller Mittel. Eine zentrale Herausforderung bestand darin, die Leistung und den Ressourceneinsatz des von mir beauftragten \gls{wsl} zu evaluieren und bei nachgewiesenen Abweichungen begründete Korrekturmaßnahmen einzuleiten, um das Projekt in Time, Scope und Budget zu halten.

% Action
\par Wie in \ref{para:04.04.05.02_OwnershipUndCommitment} geschildert, war die Leistung meines Counterparts in der Rolle des \gls{wsl} nicht zufriedenstellend, was nicht nur zu einem unnötigen Ressourcenverbrauch bei ihm führte, sondern auch zu einem erheblichen Mehrverbrauch der raren internen Ressourcen führte (vgl. \autoref{fig:CutOverPlanGU}, \autoref{fig:R2HErgebnisdokumentationStreamI2C}). Die Evaluation erfolgte anhand des CutOver-Plans sowie der Ergebnisdokumentation, die wiederholt deutliche Abweichungen zwischen geplantem und tatsächlichem Ressourceneinsatz aufzeigten. In Anbetracht der Bedeutung der Rolle des \gls{wsl} für den Erfolg des Projekts, insbesondere in Bezug auf die Einhaltung von Zeitplänen, Umfang und Budget, war es notwendig, ihn als Korrekturmaßnahme aus seiner Rolle zu entfernen und die Verantwortung für die Aufgaben, die er nicht erfüllen konnte, auf andere Teammitglieder zu verteilen (vgl. \autoref{fig:AnregungenW2undLessonsLearnedW1}).

% Result
\par Die neu zugeordneten Teammitglieder übernahmen die Verantwortung für die verbleibende Projektlaufzeit, wodurch die Ressourcen zielgerichteter eingesetzt und das Projekt planmäßig abgeschlossen werden konnte (vgl. Veränderung \autoref{fig:OrganigrammR2H}). Für Welle 1.1 ff. wurden meine Vorschläge und Anregungen vollständig von der \gls{pl} übernommen, um die Rolle des \gls{wsl} strukturell zu stärken und den Ressourceneinsatz in künftigen Wellen effizienter zu gestalten (vgl. \autoref{fig:AnregungenW2undLessonsLearnedW1}).

% Reflect
\par Die Erfahrung hat gezeigt, dass Ressourcenprobleme auf Führungsebene frühzeitig und konsequent adressiert werden müssen: Eine verzögerte Reaktion auf nachgewiesene Leistungsdefizite verursacht nicht nur direkten Ressourcenmehrverbrauch, sondern belastet auch das gesamte Projektteam. Rückblickend hätte ich die formale Neubewertung der \gls{wsl}-Rolle früher einleiten sollen, sobald die ersten Abweichungen im CutOver-Plan sichtbar wurden, anstatt auf eine Selbstkorrektur zu warten. Die daraus entwickelten Kriterien zur frühzeitigen Identifikation von Ressourcenproblemen sowie die strukturellen Anpassungen der \gls{wsl}-Rolle wurden für Welle 1.1 ff. institutionalisiert und haben die Effizienz des Ressourceneinsatzes in den Folgewellen nachweislich verbessert.